\documentclass[11pt]{article}
\usepackage[T1]{fontenc}
\usepackage[utf8]{inputenc}
\usepackage{lmodern}
\usepackage{amsmath,amssymb,amsthm,mathtools}
\usepackage[margin=1in]{geometry}
\usepackage[numbers,sort&compress]{natbib}
\usepackage{microtype}
\usepackage{hyperref}
\hypersetup{colorlinks=true,linkcolor=blue,citecolor=blue,urlcolor=blue,
 pdftitle={Positivity and recurrence formulas for oscillatory Gegenbauer Hankel determinants},
 pdfauthor={Anonymous}}
\numberwithin{equation}{section}
\newtheorem{theorem}{Theorem}[section]
\newtheorem{lemma}[theorem]{Lemma}
\newtheorem{proposition}[theorem]{Proposition}
\newtheorem{corollary}[theorem]{Corollary}
\theoremstyle{remark}
\newtheorem{remark}[theorem]{Remark}
\newcommand{\ii}{\mathrm{i}}
\newcommand{\dd}{\,\mathrm{d}}
\newcommand{\D}{D}
\newcommand{\R}{\mathbb{R}}
\newcommand{\C}{\mathbb{C}}
\newcommand{\Q}{\mathbb{Q}}
\newcommand{\N}{\mathbb{N}}
\newcommand{\wt}{\widetilde}
\DeclareMathOperator{\tr}{tr}
\DeclareMathOperator{\diag}{diag}
\DeclareMathOperator{\spanop}{span}
\allowdisplaybreaks[2]
\setlength{\parskip}{3pt}
\setlength{\emergencystretch}{2em}

\title{Positivity and recurrence formulas for\texorpdfstring{\\}{}oscillatory Gegenbauer Hankel determinants}
\author{Anonymous}
\date{Version 0.1.0-candidate, 4 November 2026. Archived at Zenodo,
DOI \href{https://doi.org/10.5281/zenodo.23933929}{10.5281/zenodo.23933929}
(all versions: \href{https://doi.org/10.5281/zenodo.23933928}{10.5281/zenodo.23933928}).}

\begin{document}
\maketitle

\begin{abstract}
For $\alpha>-1$, consider the Hankel determinants of the complex weight
$(1-t^2)^\alpha e^{\ii xt}$ on $[-1,1]$. We prove that every even-size
determinant is strictly positive for every real frequency $x$. The argument
extends the published exponent-zero theorem by continuation in $\alpha$.
Its two ingredients are a Pearson-identity proof that adjacent determinants
never vanish together, and a large-frequency expansion uniform on compact
exponent intervals. The first ingredient also implies that every complex
zero of every determinant is simple. Odd-size determinants have infinitely
many real zeros, and their first endpoint correction is computed explicitly.
The same expansion proves two formal polynomial-degree conjectures stated
in 2008 and gives complete proofs of two recurrence-coefficient conjectures
stated in 2026. An earlier 2026 preprint claimed to derive those
recurrence formulas; we show that its key determinant formula fails as
printed, and our proofs do not use it. An exact
Bessel--Jacobi determinant identity transfers the even-size result to
spherical-Bessel derivative determinants. Combined with the published
rank-three result (an unrefereed preprint), it proves the critical-length conjecture for the
eight-dimensional space spanned by $x^j\cos x,x^j\sin x$, $0\le j\le3$.
The general critical-length conjecture remains unresolved here.
\end{abstract}

\noindent\textbf{Keywords.} Hankel determinant; oscillatory orthogonality;
Gegenbauer weight; Bessel function; critical length; endpoint asymptotics.

\section{Introduction and main result}

For $\alpha>-1$ and $z\in\C$, define
\begin{equation}\label{eq:definitions}
 \mu_j^{(\alpha)}(z)=\int_{-1}^1t^j(1-t^2)^\alpha e^{\ii zt}\dd t,
 \qquad
 D_N^{(\alpha)}(z)=
 \det[\mu_{i+j}^{(\alpha)}(z)]_{i,j=0}^{N-1},\qquad D_0^{(\alpha)}=1.
\end{equation}
For fixed real $\alpha>-1$, these are entire functions of the frequency.
The determinants are real and even on the real axis. At frequency zero the measure is
positive, so every determinant is positive. At nonzero frequency the
measure is complex and ordinary positive-measure orthogonality no longer
ensures that the determinants are nonzero.

Our principal statement is the following.

\begin{theorem}\label{thm:main}
For every integer $r\ge1$, every real $\alpha>-1$, and every $x\in\R$,
\begin{equation}\label{eq:main}
                         D_{2r}^{(\alpha)}(x)>0.
\end{equation}
Consequently, the monic orthogonal polynomial of each even degree exists
uniquely for $(1-t^2)^\alpha e^{\ii xt}$ on $[-1,1]$, at every real frequency.
\end{theorem}

\paragraph{Why this matters, in plain terms.}
Integrals of rapidly oscillating functions, such as
$\int_{-1}^1f(t)(1-t^2)^\alpha e^{\ii xt}\dd t$, are expensive to compute
by ordinary rules once the frequency $x$ is large. One family of
efficient methods builds Gaussian-type rules from polynomials that are
``orthogonal'' with respect to the complex weight itself. Because that
weight is complex, such polynomials need not exist: the polynomial of
degree $N$ exists exactly when the Hankel determinant $D_N$ is nonzero.
Theorem~\ref{thm:main} says that, for every Gegenbauer exponent and every
real frequency, this never fails in even degree. It is an existence
statement only; it does not by itself show that the resulting quadrature
nodes are distinct or well conditioned. For this numerical context see
Deaño and Huybrechs \cite{DeanoHuybrechs2009}, Asheim, Deaño, Huybrechs
and Wang \cite{ADHW2014}, and the monograph \cite{DHI2017}.

\paragraph{How the paper is organised.}
Section~\ref{sec:pearson} proves that two consecutive determinants never
vanish together, so all zeros are simple. Section~\ref{sec:endpoint}
derives large-frequency expansions with constants uniform in the
exponent, reducing everything to one-variable endpoint integrals.
Section~\ref{sec:continuation} combines these to prove
Theorem~\ref{thm:main}. Sections~\ref{sec:degree}--\ref{sec:recurrence}
use the same expansions for the 2008 degree conjectures and the 2026
recurrence conjectures. Section~\ref{sec:consequences} records
consequences and limits, Section~\ref{sec:scope} the checks and the
priority record, and the appendix the Bessel--Jacobi identity. Readers
interested only in the main theorem need Sections~2--4.

This does not assert existence of the entire orthogonal-polynomial sequence
at every frequency. Each odd-size determinant has infinitely many positive
real zeros; see Corollary~\ref{cor:oddzeros}. The distinction between even
and odd degree is therefore necessary.

The case $\alpha=0$ is the even-degree existence theorem of
Celsus, Dea\~no, Huybrechs and Iserles
\cite[Theorem~4.3]{Celsus2022}. We use that published result as a base
case. The extension to all $\alpha>-1$ follows from a continuation
argument. First, a Pearson identity rules out simultaneous vanishing of
adjacent determinants. The Toda identity then rules out multiple zeros.
Second, a uniform endpoint expansion shows that an even-size determinant
is positive at every sufficiently large real frequency, with a common
threshold on each compact exponent interval. Real zeros consequently
cannot be created during continuation, enter through frequency zero, or
arrive from infinity.

The oscillatory Gegenbauer family was studied by Milovanovi\'c,
Cvetkovi\'c and Marjanovi\'c \cite{MCM2008}. Their parameter is
$\lambda=\alpha+\tfrac12$. They prove existence of the entire sequence
when $\lambda>-\tfrac12$ is rational and the frequency is a nonzero zero
of $J_{\lambda-1}$. They also conjecture degrees of two families of
polynomial Hankel expressions. Lyu and Zhou \cite{LyuZhou2026} revisit
these expressions, propose leading-coefficient products, and retain special
parameter and frequency restrictions in their existence assumptions.
Sections~\ref{sec:degree} and \ref{sec:numerator} prove both formal degree
statements and give their leading coefficients. In the shifted case,
the generic degree has two exceptional specialisations in the admissible
real parameter interval.

Section~\ref{sec:recurrence} then proves all the formulas in
\cite[Conjectures~2.5 and 2.7]{LyuZhou2026}, retaining their precise
index ranges. Those conjectures assert that the recurrence coefficients
are given by explicit rational expressions in certain monic polynomials
$X_j,Y_j$ of degree $j$, which Lyu and Zhou list only for low indices.
What is proved here is therefore: an explicit construction of such
polynomials for every index, the degree and leading-coefficient facts
that make them monic of the stated degree, and the stated identities.
Our polynomials coincide with those listed by Lyu and Zhou at the
indices they give (Section~\ref{sec:recurrence}).

\paragraph{Prior claim.} A preprint by Ashfaque \cite{Ashfaque2026},
posted on ResearchGate on 24 July 2026, claims that the Lyu--Zhou
recurrence formulas follow from a discrete Painlev\'e~IV and
tau-function identification. It therefore precedes this paper as a
public claim on Conjectures~2.5 and 2.7. As we explain in
Section~\ref{sec:scope}, the determinant formula on which that
derivation rests (its Proposition~3.1) fails as printed, so we regard
the proofs here as the first complete ones known to us for those
conjectures, not as the first claim. As far as our reading of that
preprint goes, it does not state Theorem~\ref{thm:main} or the 2008
degree results. These algebraic results use the endpoint asymptotics and
determinant identities; they do not require Theorem~\ref{thm:main}.
The positivity theorem additionally uses the exponent-zero base case
and continuation.

The exact Fourier--Gegenbauer Wronskian object also occurs in
Dimitrov and Xu \cite{DimitrovXu2019}. Their positive-definiteness results
for Fourier Wronskians are not pointwise positivity statements. Theorem
\ref{thm:main} concerns pointwise positivity at every even size. An exact
atomic example in Section~\ref{sec:consequences} illustrates why positive
definiteness, or symmetry of a positive underlying measure, does not
itself prove this conclusion.

Jacobi-type kissing polynomials are also studied by Barhoumi
\cite{Barhoumi2021} in a regime where degree and frequency grow together.
Those large-degree asymptotic existence results have a different scope
from the all-finite-even-degree, all-real-frequency assertion above.

An appendix supplies an exact identity with signed derivative Hankel
determinants of spherical-Bessel functions. This identity is derived from
the standard unitary orbital integral and the Jacobi law for random
projection compressions; no separate novelty claim is made for that
classical representation. It gives all even-size nonvanishing conditions
in the critical-length problem of Carnicer, Mainar and Pe\~na
\cite{CMP2023}. Combined with Kounchev and Render
\cite[Corollary~21]{KounchevRender2025}, it settles the case $n=3$ of that
problem; this corollary is conditional on their result, which is an
unrefereed preprint. Higher odd-size first-zero comparisons are not proved here.

The priority discussion in Section~\ref{sec:scope} records the bounded
literature comparison, including an earlier competing claim. The
mathematical statements below do not depend on treating a search result
as proof of historical priority.

\section{A Pearson identity and simple zeros}\label{sec:pearson}

Fix a real $\alpha>-1$. For a polynomial $p$, write
\[
 L_z(p)=\int_{-1}^1p(t)(1-t^2)^\alpha e^{\ii zt}\dd t.
\]
Integrating the derivative of
$R(t)(1-t^2)^{\alpha+1}e^{\ii zt}$ gives
\begin{equation}\label{eq:pearson}
 L_z\!\left((1-t^2)R'
       +[\ii z(1-t^2)-2(\alpha+1)t]R\right)=0.
\end{equation}
The endpoint values vanish because $\alpha+1>0$, and all terms are
integrable. This identity holds for every polynomial $R$ and every
complex frequency.

\begin{lemma}\label{lem:adjacent}
There is no $z\in\C\setminus\{0\}$ and no integer $n\ge1$ for which
\[
                 D_n^{(\alpha)}(z)=D_{n+1}^{(\alpha)}(z)=0.
\]
\end{lemma}

\begin{proof}
Fix a frequency and suppress $\alpha,z$ from the notation. If an adjacent
pair vanishes, choose the least such $n$. Then $D_{n-1}\ne0$; put
$d=n-1$. Nonsingularity of the size-$d$ moment matrix provides a unique
monic polynomial $p$ of degree $d$ satisfying
\begin{equation}\label{eq:orthstart}
                         L(pq)=0\qquad(\deg q\le d-1).
\end{equation}
For $d=0$, this means $p=1$ and the condition is empty. The monic
triangular basis change to $1,t,\ldots,t^{d-1},p$ gives
$D_{d+1}=D_dL(p^2)$. Hence $L(p^2)=0$, and $p$ is orthogonal to every
polynomial of degree at most $d$.

In the monic basis $1,t,\ldots,t^{d-1},p,tp$, the size-$(d+2)$ Gram
matrix has the form
\[
 \begin{pmatrix}
  H_d&0&v\\
  0&0&c\\
  v^{T}&c&e
 \end{pmatrix},\qquad c=L(tp^2).
\]
The Schur complement of $H_d$ therefore gives
$D_{d+2}=-D_dc^2$. Since $D_{d+2}=0$ and $D_d\ne0$, we have $c=0$.
The polynomial $tp$ is monic of degree $d+1$, so $p$ is orthogonal to
all polynomials of degree at most $d+1$.

Suppose inductively that this orthogonality holds through degree
$M\ge d+1$. Set $k=M-d-1\ge0$ and substitute $R=t^kp^2$ in
\eqref{eq:pearson}. Every term except
$-\ii zL(t^{k+2}p^2)$ has the form $L(pq)$ with $\deg q\le M$.
For example, the largest degree in the term involving $p'$ is
$\deg(t^{k+2}p')=k+2+d-1=M$. Since $z\ne0$, it follows that
$L(t^{k+2}p^2)=0$. The polynomial $t^{k+2}p$ is monic of degree
$M+1$, which extends the orthogonality by one degree. Induction shows
that $L(pq)=0$ for every polynomial $q$.

Approximate the continuous function $\overline{p(t)}e^{-\ii zt}$
uniformly on $[-1,1]$ by complex polynomials. Passage to the limit gives
\[
                       \int_{-1}^1|p(t)|^2(1-t^2)^\alpha\dd t=0,
\]
contradicting the fact that $p$ is monic.
\end{proof}

\begin{remark}\label{rem:jacobi}
The same proof applies to the nonsymmetric Jacobi weight
$(1-t)^a(1+t)^b e^{\ii zt}$, $a,b>-1$. Its Pearson bracket is
$\ii z(1-t^2)+(b-a)-(a+b+2)t$. The extra constant term does not change
the degree induction. The positivity theorem uses symmetry and is
asserted only for $a=b$.
\end{remark}

Because $(\mu_j^{(\alpha)})'=\ii\mu_{j+1}^{(\alpha)}$, the
Desnanot--Jacobi identity for consecutive derivative Hankel determinants
gives
\begin{equation}\label{eq:toda}
 D_N D_N''-(D_N')^2=-D_{N-1}D_{N+1},\qquad N\ge1.
\end{equation}
Here and below frequency derivatives are taken with $\alpha$ fixed.
The sign can already be read at $N=1$:
$\mu_0\mu_0''-(\mu_0')^2=-\mu_0\mu_2+\mu_1^2=-D_2$.

\begin{corollary}\label{cor:simple}
Every complex zero of every $D_N^{(\alpha)}$, $N\ge1$, is simple.
\end{corollary}

\begin{proof}
If $D_N=D_N'=0$ at a nonzero frequency, \eqref{eq:toda} forces one of
the adjacent determinants to vanish, contrary to Lemma~\ref{lem:adjacent}.
Frequency zero is not a zero, because its moment matrix is positive
definite.
\end{proof}

For later normalisation, the zero-frequency value is
\begin{equation}\label{eq:selbergzero}
 D_N^{(\alpha)}(0)=2^{N(2\alpha+N)}
 \prod_{j=0}^{N-1}
 \frac{j!\,\Gamma(\alpha+j+1)^2}
      {\Gamma(2\alpha+N+j+1)}>0.
\end{equation}
This is the Selberg integral after $t=2s-1$, or equivalently the product
of the monic Jacobi polynomial norms. Positivity alone in the preceding
argument does not require this evaluation.

\section{Endpoint expansions, uniform in the exponent}\label{sec:endpoint}

Let
\begin{equation}\label{eq:Laguerreproduct}
 L_m(\alpha)=\prod_{j=0}^{m-1}j!\,\Gamma(\alpha+j+1),\qquad L_0=1.
\end{equation}
Andr\'eief's identity yields
\begin{equation}\label{eq:andreief}
 D_N^{(\alpha)}(x)=\frac1{N!}
 \int_{[-1,1]^N}\Delta(t)^2\prod_{j=1}^N(1-t_j^2)^\alpha
                    e^{\ii x\sum_jt_j}\dd t_1\cdots\dd t_N,
\end{equation}
where $\Delta(t)=\prod_{i<j}(t_j-t_i)$.

The whole large-frequency analysis rests on one elementary fact about
one-variable endpoint integrals, which we prove first with constants that
are uniform in the exponent. Throughout, $K\subset(-1,\infty)$ is a compact
interval, and $O_K(\cdot)$ denotes a bound whose constant depends on $K$
and on the displayed fixed integers, but not on $\alpha\in K$ or on $x$.

\begin{lemma}[uniform endpoint lemma]\label{lem:watson}
Fix $\delta\in(0,1]$, a smooth cutoff $\chi$ on $[0,\infty)$ with
$\chi=1$ on $[0,\delta/3]$ and $\chi=0$ on $[2\delta/3,\infty)$, an
integer $q\ge0$, and a function $h(s,\alpha)$, smooth in
$s\in[0,\delta]$, whose $s$-derivatives of every order are bounded and
jointly continuous on $[0,\delta]\times K$. Then for every integer $M\ge1$
and $x\ge1$,
\begin{equation}\label{eq:watson}
 \int_0^\delta s^{\alpha+q}h(s,\alpha)\chi(s)e^{-\ii xs}\dd s
 =\sum_{m=0}^{M-1}\frac{\partial_s^m h(0,\alpha)}{m!}\,
   \Gamma(\alpha+q+m+1)(\ii x)^{-(\alpha+q+m+1)}
 +O_K\bigl(x^{-(q+M)}\bigr),
\end{equation}
with the principal branch of $(\ii x)^{-c}=x^{-c}e^{-\ii\pi c/2}$.
The coefficients are continuous in $\alpha\in K$. The same holds with
$e^{+\ii xs}$ and $(-\ii x)$ in place of $e^{-\ii xs}$ and $(\ii x)$.
\end{lemma}

\begin{proof}
Taylor's theorem gives
$h(s,\alpha)=\sum_{m<M}h_m(\alpha)s^m+s^MR(s,\alpha)$ with
$h_m=\partial_s^mh(0,\alpha)/m!$, where every $s$-derivative of $R$ is
bounded on $[0,\delta]\times K$. Put $n=q+M$ and
$g(s)=s^{\alpha+q+M}R(s,\alpha)\chi(s)$. For $k\le n$ the derivative
$g^{(k)}$ is a finite sum of terms bounded by a constant times
$s^{\alpha+q+M-k}$, and $\alpha+q+M-k\ge\alpha>-1$. Hence every
$g^{(k)}$, $k\le n$, is integrable on $[0,\delta]$ with
$\int_0^\delta|g^{(k)}|\le C_K$, since
$\int_0^\delta s^{\alpha}\dd s\le\delta^{\alpha+1}/(\alpha+1)$ is
bounded on $K$. For $k<n$, $g^{(k)}(s)=O(s^{\alpha+1})\to0$ as
$s\to0$, and $g^{(k)}(\delta)=0$. Integrating by parts $n$ times
therefore gives $\bigl|\int_0^\delta g\,e^{-\ii xs}\dd s\bigr|\le
C_Kx^{-n}$.

It remains to treat $\int_0^\infty s^\beta\chi(s)e^{-\ii xs}\dd s$ with
$\beta=\alpha+q+m$. Insert the factor $e^{-\varepsilon s}$,
$\varepsilon>0$. Then
\[
 \int_0^\infty s^\beta e^{-(\varepsilon+\ii x)s}\dd s
 =\Gamma(\beta+1)(\varepsilon+\ii x)^{-\beta-1}
\]
and the difference from the cutoff integral is
$\int_0^\infty s^\beta(1-\chi(s))e^{-(\varepsilon+\ii x)s}\dd s$.
The factor $1-\chi$ vanishes near zero. Choose one integer
$L\ge q+M$ with $L>\max_{\alpha\in K,\,m<M}(\alpha+q+m+1)$, and
integrate by parts $L$ times; for $\varepsilon>0$ all boundary terms
vanish. By the Leibniz rule the $L$th derivative of
$s^\beta(1-\chi(s))e^{-\varepsilon s}$ is a sum of compactly supported
cutoff-derivative terms and terms
$\varepsilon^js^{\beta-L+j}e^{-\varepsilon s}$ on $s\ge\delta/3$, and
$\varepsilon^js^je^{-\varepsilon s}\le C_j$. Hence it is dominated,
uniformly in $\varepsilon\in(0,1]$ and $\alpha\in K$, by an integrable
multiple of $s^{\beta-L}$ on $[\delta/3,\infty)$. This difference is therefore
$O_K(x^{-L})$ uniformly in $\varepsilon$. Letting $\varepsilon\to0$ by
dominated convergence on the cutoff side gives
$\int_0^\infty s^\beta\chi e^{-\ii xs}\dd s
=\Gamma(\beta+1)(\ii x)^{-\beta-1}+O_K(x^{-L})$, with $L$ as large as
required. Summing over $m<M$ proves \eqref{eq:watson}. Continuity of the
coefficients follows from continuity of $\Gamma$ and of $h_m$. When
$\alpha>0$ the remainder bound is weaker than the last retained term;
this is harmless because $M$ may be taken as large as required.
\end{proof}

\begin{lemma}[moment splitting]\label{lem:momentsplit}
Fix $\delta\in(0,1/2]$ and smooth functions $\chi_+,\chi_0,\chi_-$ on
$[-1,1]$, with values in $[0,1]$ and sum one, where $\chi_+(t)=\chi(1-t)$,
$\chi_-(t)=\chi(1+t)$, and $\chi_0$ is supported in
$[-1+\delta/3,1-\delta/3]$. For each $j\ge0$ write
$\mu_j^{(\alpha)}=\mu_j^++\mu_j^0+\mu_j^-$ according to this partition. Then,
for real $x\ge1$, $\mu_j^0(x)=O_K(x^{-L})$ for every $L$, while
$e^{-\ii x}\mu_j^+(x)$ and $e^{\ii x}\mu_j^-(x)$ have expansions of the form
\eqref{eq:watson} to every order, with coefficients continuous in
$\alpha\in K$.
\end{lemma}

\begin{proof}
On the support of $\chi_0$ the function $t^j(1-t^2)^\alpha\chi_0(t)$ is
smooth, compactly supported in $(-1,1)$, and has all derivatives
bounded on $K$; repeated integration by parts gives the first claim.
Near $t=1$, put $t=1-s$. Then
$t^j(1-t^2)^\alpha e^{\ii xt}
 =e^{\ii x}s^\alpha\,(1-s)^j(2-s)^\alpha e^{-\ii xs}$, and
$h(s,\alpha)=(1-s)^j(2-s)^\alpha$ satisfies the hypotheses of
Lemma~\ref{lem:watson} on $[0,\delta]$, because $2-s\ge3/2$ there.
The endpoint $t=-1$ is symmetric.
\end{proof}

\begin{proposition}\label{prop:corners}
Fix $N\ge1$ and a compact interval $K\subset(-1,\infty)$. For real
$x\ge1$, $D_N^{(\alpha)}(x)$ is a sum over endpoint allocations
$k+l=N$ of terms $e^{\ii(k-l)x}F_{k,l}(x,\alpha)$ plus $O_K(x^{-L})$
for any prescribed $L$. Each $F_{k,l}$ has an asymptotic expansion in
powers $x^{-N(\alpha+1)-m}$, $m=0,1,2,\ldots$, to every order, with
coefficients continuous in $\alpha\in K$ and remainders uniform for
$\alpha\in K$. The first nonzero term of $e^{\ii(k-l)x}F_{k,l}$ is
\begin{equation}\label{eq:corner}
 2^{N\alpha+2kl}L_k(\alpha)L_l(\alpha)
 x^{-E_k}
 \exp\!\left(\ii(k-l)
             \left[x-\frac\pi2(\alpha+N)\right]\right),
 \quad E_k=N\alpha+k^2+l^2.
\end{equation}
The frequency derivative $\partial_xD_N^{(\alpha)}$ has the
termwise-differentiated expansion, with the same uniformity.
\end{proposition}

\begin{proof}
\emph{Reduction to one-variable integrals.} Expanding $\Delta(t)^2$ into
monomials in \eqref{eq:andreief} expresses $D_N^{(\alpha)}$ as a fixed finite
linear combination of products $\prod_{i=1}^N\mu_{a_i}^{(\alpha)}$
(this is just the Leibniz expansion of the Hankel determinant). Insert
the splitting of Lemma~\ref{lem:momentsplit} into each factor. Any
product containing a factor $\mu^0$ is $O_K(x^{-L})$, since every
$|\mu_j^\pm|$ is bounded by $\int_{-1}^1(1-t^2)^\alpha\dd t$, which is
bounded on $K$. Collect the remaining products according to how many
factors carry $+$ and how many carry $-$. Equivalently, for a fixed
allocation in which $k$ coordinates carry $\chi_+$ and $l$ carry
$\chi_-$, the corresponding term of \eqref{eq:andreief} is
\[
 \binom Nk\frac1{N!}\int\Delta(t)^2
   \prod_{i\le k}(1-t_i^2)^\alpha\chi_+(t_i)e^{\ii xt_i}
   \prod_{i>k}(1-t_i^2)^\alpha\chi_-(t_i)e^{\ii xt_i}\dd t .
\]
Put $t_i=1-s_i$ for $i\le k$ and $t_i=-1+s_i'$ for $i>k$. The square
$\Delta(t)^2$ becomes the \emph{polynomial}
$\Delta(s)^2\Delta(s')^2\prod_{i\le k<j}(2-s_i-s_j')^2$. All other
factors are products of one-variable functions of the type in
Lemma~\ref{lem:momentsplit}. Expanding this polynomial into monomials
therefore writes the allocation term exactly as
$e^{\ii(k-l)x}F_{k,l}$, where $F_{k,l}$ is a finite sum of products of
one-variable integrals covered by Lemma~\ref{lem:watson}. No
multivariable Taylor remainder occurs, and the uniformity and
continuity statements follow by multiplying finitely many uniform
one-variable expansions.

\emph{The leading term.} A monomial of total degree $d$ in the
$s,s'$ variables contributes, by \eqref{eq:watson}, terms of order
$x^{-N(\alpha+1)-d-m}$, $m\ge0$. The polynomial
$\Delta(s)^2\Delta(s')^2$ is homogeneous of degree
$k(k-1)+l(l-1)$ and is the lowest-degree part of the full polynomial;
its coefficient multiplies $2^{2kl}$ from the cross factor at
$s=s'=0$. The constant terms of the smooth factors are
$h(0,\alpha)=2^\alpha$ per coordinate, giving $2^{N\alpha}$. The order is
therefore $N(\alpha+1)+k(k-1)+l(l-1)=E_k$, and the leading
coefficient is obtained by replacing each one-variable integral by its
first term in \eqref{eq:watson}. That replacement is the same as
evaluating
\begin{equation}\label{eq:Laguerreintegral}
 \int_{[0,\infty)^m}\Delta(u)^2\prod_{j=1}^m u_j^\alpha
                 e^{-s\sum_j u_j}\dd u
 =m!L_m(\alpha)s^{-m(\alpha+m)},\qquad \Re s>0,
\end{equation}
at the boundary values $s=\pm\ii x$, where the right-hand side is read
as the analytic continuation from $\Re s>0$. Indeed, both sides of
\eqref{eq:Laguerreintegral} are sums of the same products of Gamma
values times powers of $s$, and \eqref{eq:watson} supplies exactly these
products with $s=\ii x$ (respectively $-\ii x$). At $s=1$, Andr\'eief
reduces the integral to $m!\det[\Gamma(\alpha+i+j+1)]$. Factoring
$\Gamma(\alpha+i+1)$ from row $i$ leaves $(\alpha+i+1)_j$, a monic
polynomial of degree $j$ in $i$. Its determinant is the Vandermonde
determinant at $0,\ldots,m-1$, namely $\prod_{j=0}^{m-1}j!$. Homogeneity
gives the formula for general $s$. The boundary values supply the
phases $e^{\mp\ii\pi m(\alpha+m)/2}$. The factor $\binom Nk$ and $1/N!$
cancel the two factorials in \eqref{eq:Laguerreintegral}. The phase
simplifies using $k(\alpha+k)-l(\alpha+l)=(k-l)(\alpha+N)$, giving
\eqref{eq:corner}. This coefficient is nonzero, so it is the first
nonzero term.

\emph{The derivative.} For the one-variable integrals
$I_q(x)=\int s^{\alpha+q}h\chi e^{-\ii xs}\dd s$, differentiation under the
integral gives $I_q'=-\ii I_{q+1}$. Applying Lemma~\ref{lem:watson} to
both sides and using $\Gamma(p+1)=p\Gamma(p)$ shows that the expansion of
$I_q'$ is the termwise derivative of that of $I_q$; the conjugate phase
is analogous. The product rule, applied to each allocation term
including its phase $e^{\ii(k-l)x}$, then gives the uniform expansion of
$\partial_xD_N^{(\alpha)}$ as the termwise derivative, without
differentiating an uncontrolled remainder.

\emph{Joint continuity.} By dominated convergence, every
$\mu_j^{(\alpha)}(x)$ and its $x$-derivatives are jointly continuous in
$(\alpha,x)\in(-1,\infty)\times\R$; hence so are $D_N^{(\alpha)}(x)$ and
$\partial_xD_N^{(\alpha)}(x)$. This is used in Section~\ref{sec:continuation}.
\end{proof}

\begin{corollary}\label{cor:eventail}
For $r\ge1$, uniformly on compact $K\subset(-1,\infty)$,
\begin{equation}\label{eq:eventail}
 D_{2r}^{(\alpha)}(x)=C_{r,\alpha}
 x^{-2r(\alpha+r)}\bigl(1+O_K(x^{-2})\bigr),
 \qquad
 C_{r,\alpha}=2^{2r(\alpha+r)}L_r(\alpha)^2>0.
\end{equation}
In particular, there is a common $R$ such that
$D_{2r}^{(\alpha)}(x)>0$ for every $\alpha\in K$ and $x\ge R$.
\end{corollary}

\begin{proof}
Write $k=r+s$, $l=r-s$ in \eqref{eq:corner}. Then
$E_k=2r(\alpha+r)+2s^2$. The balanced allocation $k=l=r$ is the
unique leading term, and its phase is one. Its first correction
cancels between the conjugate endpoint groups; the explicit calculation
below also verifies this cancellation. Every unbalanced allocation is
smaller by at least $x^{-2}$. The coefficient is continuous and strictly
positive on $K$, so the uniform expansion implies a common positivity
threshold.
\end{proof}

The odd case needs a first correction for the degree calculations later.

\begin{proposition}\label{prop:oddexpansion}
Let $r\ge0$, set $N=2r+1$, and put
\begin{align}
 E&=(2r+1)\alpha+2r^2+2r+1,&
 \theta&=x-\frac{\pi\alpha}{2},\label{eq:oddparameters}\\
 A_{r,\alpha}&=2^{(2r+1)\alpha+2r(r+1)+1}
                  L_r(\alpha)L_{r+1}(\alpha),&
 T_{r,\alpha}&=\frac{\alpha(\alpha+2r+1)}2+r(r+1).
 \label{eq:oddconstants}
\end{align}
Then, uniformly on compact exponent intervals,
\begin{equation}\label{eq:oddexpansion}
 D_{2r+1}^{(\alpha)}(x)=(-1)^r A_{r,\alpha}x^{-E}
 \left[\sin\theta+\frac{T_{r,\alpha}}x\cos\theta+O_K(x^{-2})\right].
\end{equation}
The bracket remainder can be differentiated once with error
$O_K(x^{-2})$.
\end{proposition}

\begin{proof}
The leading allocations are $(k,l)=(r+1,r)$ and $(r,r+1)$.
Their phases in \eqref{eq:corner} are conjugate and add to
$2(-1)^r\sin\theta$. The next allocations are down by $x^{-4}$.
In the scaled endpoint variables, the linear correction to the smooth
amplitude at an allocation $(k,l)$ is
\begin{equation}\label{eq:linearcorner}
 -\frac1x\left[
    \left(\frac\alpha2+l\right)\sum_{j=1}^k u_j
   +\left(\frac\alpha2+k\right)\sum_{j=1}^l v_j\right].
\end{equation}
Differentiate \eqref{eq:Laguerreintegral} in $s$. For the group with
$e^{-\ii\sum u_j}$, the normalised endpoint moment of $\sum u_j$ is
$-\ii k(\alpha+k)$; the conjugate group has moment
$\ii l(\alpha+l)$. For $(k,l)=(r+1,r)$, the relative correction
is $\ii T_{r,\alpha}/x$, where
\[
 T_{r,\alpha}=
 \left(\frac\alpha2+r\right)(r+1)(\alpha+r+1)
 -\left(\frac\alpha2+r+1\right)r(\alpha+r).
\]
Expansion gives \eqref{eq:oddconstants}. Adding the conjugate allocation
proves \eqref{eq:oddexpansion}. The remainder and its differentiated
version follow from Proposition~\ref{prop:corners}. For the balanced
even allocation, the two terms in \eqref{eq:linearcorner} instead
cancel, as used in Corollary~\ref{cor:eventail}.
\end{proof}

\section{Continuation and the real zeros}\label{sec:continuation}

\begin{proof}[Proof of Theorem~\ref{thm:main}]
For real frequency, reflection of the integration variable gives
$\overline{\mu_j^{(\alpha)}(x)}=(-1)^j\mu_j^{(\alpha)}(x)$.
The row and column sign exponents of a size-$N$ determinant sum to
$N(N-1)$, an even number. Thus $D_N^{(\alpha)}(x)$ is real; the same
calculation shows it is even in $x$.

Fix $2r$ and a target exponent $\alpha_*>-1$. Let $K$ be the closed
interval joining $0$ and $\alpha_*$. Positive definiteness at frequency
zero, compactness, and continuity give $\varepsilon>0$ such that
$D_{2r}^{(\alpha)}(x)>0$ for all $\alpha\in K$ and $|x|\le\varepsilon$.
Corollary~\ref{cor:eventail} gives $R>\varepsilon$ with the same
positivity for $|x|\ge R$.

Define
\[
 S=\{\alpha\in K:D_{2r}^{(\alpha)}(x)>0
                       \text{ for every }x\in[\varepsilon,R]\}.
\]
The exponent-zero theorem \cite[Theorem~4.3]{Celsus2022} states that
for every $n\ge0$ and every real frequency the monic kissing polynomial
of degree $2n$ for $e^{\ii\omega t}$ on $[-1,1]$ exists; it is
unconditional, with no asymptotic restriction. It gives
$0\in S$. Indeed, it proves nonvanishing, and positivity follows by
continuity from $x=0$. The set $S$ is relatively open, since a positive
continuous function on this compact frequency interval has a positive
minimum.

It is also relatively closed. If $\alpha_j\in S$ tends to $\alpha\in K$,
then $D_{2r}^{(\alpha)}\ge0$ on $[\varepsilon,R]$. A zero would lie
strictly inside the interval because its endpoints are already in
positive regions. At such an interior minimum the frequency derivative
would also vanish, contradicting Corollary~\ref{cor:simple}. Hence
$\alpha\in S$. Connectedness gives $S=K$, and evenness completes the
proof for every real frequency.
\end{proof}

\begin{corollary}\label{cor:oddzeros}
For every real $\alpha>-1$ and every $r\ge0$, the function
$D_{2r+1}^{(\alpha)}$ has infinitely many positive real zeros, all simple.
For all sufficiently large integers $m$, exactly one zero lies near
$x_m=m\pi+\pi\alpha/2$, and it equals
\begin{equation}\label{eq:oddzerolocation}
                      x_m-\frac{T_{r,\alpha}}{x_m}+O(m^{-2}).
\end{equation}
This local label need not be the global index of that zero.
\end{corollary}

\begin{proof}
The leading sine in \eqref{eq:oddexpansion} has alternating signs at its
successive extrema. The differentiated expansion gives a derivative
of fixed nonzero sign in a small neighbourhood of each sufficiently
large phase zero, and the function is bounded away from zero outside
these neighbourhoods. This proves existence and uniqueness there.
Substituting $x=x_m+\delta$ in \eqref{eq:oddexpansion} gives
$\delta=-T_{r,\alpha}/x_m+O(m^{-2})$. Simplicity of every zero,
including the remaining ones, is Corollary~\ref{cor:simple}.
\end{proof}

The Toda identity and Theorem~\ref{thm:main} also give a direct
consequence between successive odd-size zeros:
\begin{equation}\label{eq:logconcavity}
 \frac{\mathrm{d}^2}{\mathrm{d}x^2}
          \log|D_{2r+1}^{(\alpha)}(x)|
 =-\frac{D_{2r}^{(\alpha)}(x)D_{2r+2}^{(\alpha)}(x)}
          {D_{2r+1}^{(\alpha)}(x)^2}<0.
\end{equation}
This strict logarithmic concavity does not by itself locate the first
positive zero.

\section{The degree of the Hankel polynomials}\label{sec:degree}

Let $u$ be a polynomial variable. For a parameter $\lambda$, set
$P_{-1}=0$, $P_0=1$, $P_1=-2\lambda$, and
\begin{equation}\label{eq:Prec}
 P_{k+2}=-(k+2\lambda+1)P_{k+1}-uP_k-kuP_{k-1},\qquad k\ge0.
\end{equation}
Define
\begin{equation}\label{eq:Hdef}
 H_n(u;\lambda)=\det[P_{i+j}(u;\lambda)]_{i,j=0}^{n-1},\qquad H_0=1.
\end{equation}
These are precisely the polynomials in \cite[pp.~50--53]{MCM2008},
with $u=\zeta^2$. Put
\begin{equation}\label{eq:Kdef}
 a=\lambda+\frac12,\qquad
 K_r(a)=\prod_{j=0}^{r-1}j!(a)_j,\qquad K_0=1.
\end{equation}

\begin{theorem}\label{thm:Hdegree}
For every real $\lambda>-\tfrac12$ and every $n\ge0$,
\begin{equation}\label{eq:Hdegree}
                       \deg_u H_n=\left\lfloor\frac{n^2}{4}\right\rfloor.
\end{equation}
The leading coefficients $\ell_n(\lambda)$ are
\begin{align}
 \ell_{2r}(\lambda)&=(-1)^r2^{2r(r-1)}K_r(a)^2,
                    &&r\ge1,\label{eq:Hevenlead}\\
 \ell_{2r+1}(\lambda)&=2^{2r^2}K_r(a)K_{r+1}(a),
                    &&r\ge0.\label{eq:Hoddlead}
\end{align}
The degree upper bound and these coefficient formulas hold as
polynomial identities for every complex $\lambda$. Equality in the
degree holds whenever the displayed coefficient is nonzero.
\end{theorem}

\begin{proof}
First assume $\lambda>-\tfrac12$ is real and set
$\alpha=\lambda-\tfrac12$. The Fourier--Bessel integral is
\begin{equation}\label{eq:muBessel}
 \mu_0^{(\alpha)}(\zeta)=
 \sqrt\pi\,\Gamma\left(\lambda+\frac12\right)
                    \left(\frac2\zeta\right)^\lambda J_\lambda(\zeta).
\end{equation}
Let $\zeta$ run through an unbounded sequence of positive zeros of
$J_{\lambda-1}$. At these frequencies the Bessel recurrence gives
$\mu_1=2\ii\lambda\mu_0/\zeta$. The Pearson identity for a monomial
gives
\[
 \ii\zeta\mu_{k+2}
   =k\mu_{k-1}-(k+2\lambda+1)\mu_{k+1}+\ii\zeta\mu_k.
\]
Comparison with \eqref{eq:Prec} proves
\begin{equation}\label{eq:Pmoment}
 \mu_j^{(\alpha)}(\zeta)
 =\mu_0^{(\alpha)}(\zeta)(\ii\zeta)^{-j}P_j(\zeta^2;\lambda).
\end{equation}
Factoring rows and columns gives the exact identity
\begin{equation}\label{eq:DH}
 D_n^{(\alpha)}(\zeta)=
 \mu_0^{(\alpha)}(\zeta)^n(\ii\zeta)^{-n(n-1)}H_n(\zeta^2;\lambda).
\end{equation}
This step assumes neither existence of an orthogonal-polynomial
sequence nor rationality of $\lambda$.

Write $\theta=\zeta-\pi\alpha/2$ and
$B=2^{\alpha+1}\Gamma(\alpha+1)$. The large-argument Bessel expansion
at zeros of $J_{\lambda-1}$ gives
\begin{equation}\label{eq:phasezeros}
 \cos\theta=O(\zeta^{-1}),\qquad
 \sin^2\theta=1+O(\zeta^{-2}),\qquad
 \mu_0^{(\alpha)}(\zeta)
 =B\zeta^{-\alpha-1}\bigl(\sin\theta+O(\zeta^{-1})\bigr).
\end{equation}
In particular $\mu_0$ is nonzero along all sufficiently large members
of the sequence.

For $n=2r$, insert \eqref{eq:eventail} and \eqref{eq:phasezeros} into
\eqref{eq:DH}. Since $\ii^{n(n-1)}=(-1)^r$ and
$(\sin\theta)^{2r}\to1$,
\[
 \frac{H_{2r}(\zeta^2;\lambda)}{\zeta^{2r^2}}
 \longrightarrow
 (-1)^r\frac{2^{2r(\alpha+r)}L_r(\alpha)^2}{B^{2r}}
 =(-1)^r2^{2r(r-1)}K_r(a)^2.
\]
For $n=2r+1$, the leading odd term is
$(-1)^rA_{r,\alpha}\zeta^{-E}\sin\theta$.
Again $\ii^{n(n-1)}=(-1)^r$, so the two signs cancel and
\[
 \frac{H_{2r+1}(\zeta^2;\lambda)}{\zeta^{2r(r+1)}}
 \longrightarrow
 \frac{A_{r,\alpha}}{B^{2r+1}}
 =2^{2r^2}K_r(a)K_{r+1}(a).
\]
Each limit is finite and nonzero because $a>0$.
Since $H_n$ is a polynomial in $u$, its growth along the unbounded
sequence $u=\zeta^2$ determines its degree and leading coefficient:
a higher degree would give an unbounded normalised sequence, and a
lower degree would give limit zero.

Finally, every coefficient of $H_n$ is a polynomial in $\lambda$.
The coefficients above the stated degree vanish on the real interval
$(-\tfrac12,\infty)$ and therefore vanish identically. The leading
coefficient identities extend in the same way.
\end{proof}

Rearranging the products in \eqref{eq:Hevenlead}--\eqref{eq:Hoddlead}
gives the equivalent formulas
\begin{align}
 \ell_{2r}&=(-1)^r2^{r(r-1)}
   \prod_{j=1}^{r-1}[j(2\lambda+2j-1)]^{2(r-j)},\label{eq:LYeven}\\
 \ell_{2r+1}&=2^{r^2}
   \prod_{j=1}^{r}[j(2\lambda+2j-1)]^{2(r-j)+1}.\label{eq:LYodd}
\end{align}
These are the products proposed in \cite[Remark~2.6]{LyuZhou2026}.
The degree statement is the first conjecture on p.~53 of
\cite{MCM2008}.

\section{The shifted-numerator degree}\label{sec:numerator}

For $n\ge1$, define the shifted-last-column determinant
\begin{equation}\label{eq:Htilde}
 \wt H_n=\det[P_{i+j+\mathbf{1}_{\{j=n-1\}}}]_{i,j=0}^{n-1},
 \qquad
 V_n=\wt H_nH_{n+1}-\wt H_{n+1}H_n.
\end{equation}
Thus the last column uses $P_n,\ldots,P_{2n-1}$; all earlier columns
are unchanged. The prime used for this expression in \cite{MCM2008}
denotes a shifted determinant, not a parameter derivative.

\begin{theorem}\label{thm:Vdegree}
As a formal polynomial over $\Q(\lambda)$,
\begin{equation}\label{eq:Vdegree}
                            \deg_u V_n=\frac{n(n+1)}2.
\end{equation}
The coefficient of $u^{n(n+1)/2}$ is
\begin{equation}\label{eq:Vlead}
                       C_n(\lambda)\ell_n(\lambda)\ell_{n+1}(\lambda),
\end{equation}
where
\begin{align}
 C_{2r}(\lambda)&=(r+2)\lambda+\frac{r(2r+5)}2,
                      &&r\ge1,\label{eq:Ceven}\\
 C_{2r+1}(\lambda)&=1+\frac{3r}{2}-r^2-r\lambda,
                      &&r\ge0.\label{eq:Codd}
\end{align}
For real $\lambda>-\tfrac12$, the specialised degree equals
$n(n+1)/2$ except at $(n,\lambda)=(3,3/2)$ and $(5,0)$, where it is
respectively $5$ and $14$.
\end{theorem}

\begin{proof}
Define $\wt D_n$ by the same shifted-last-column construction with the
moments $\mu_j^{(\alpha)}$. Differentiate $D_n$ by columns. Each
column derivative other than the last duplicates the next column, so
\begin{equation}\label{eq:shiftderivative}
                              D_n'(x)=\ii\wt D_n(x).
\end{equation}
This holds at every frequency, before any specialisation. At a positive
zero $\zeta$ of $J_{\lambda-1}$, \eqref{eq:Pmoment} gives
\[
 \wt D_n(\zeta)=\mu_0(\zeta)^n
          (\ii\zeta)^{-n(n-1)-1}\wt H_n(\zeta^2;\lambda).
\]
Consequently,
\begin{equation}\label{eq:Vlogder}
 \frac{V_n}{H_nH_{n+1}}
 =\zeta\left(\frac{D_n'(\zeta)}{D_n(\zeta)}
             -\frac{D_{n+1}'(\zeta)}{D_{n+1}(\zeta)}\right).
\end{equation}
No identity valid only on the discrete Bessel-zero set has been
differentiated.

Put $\alpha=\lambda-\tfrac12$. The even expansion and its derivative
give
\begin{equation}\label{eq:evenloglimit}
 \lim_{x\to\infty}x\frac{(D_{2r}^{(\alpha)})'(x)}{D_{2r}^{(\alpha)}(x)}
                        =-2r(\alpha+r).
\end{equation}
For the odd limit the first correction in
Proposition~\ref{prop:oddexpansion} matters. The first two Bessel terms
are
\[
 J_{\alpha-1/2}(x)=\sqrt{\frac2{\pi x}}
 \left[\cos\theta-\frac{B_\alpha}{x}\sin\theta+O(x^{-2})\right],
 \qquad B_\alpha=\frac{\alpha(\alpha-1)}2.
\]
Thus, at its large positive zeros, $\zeta\cot\theta\to B_\alpha$.
Taking the logarithmic derivative of \eqref{eq:oddexpansion} yields
\begin{align}
 \lim_{\zeta\to\infty}\zeta
 \frac{(D_{2r+1}^{(\alpha)})'(\zeta)}{D_{2r+1}^{(\alpha)}(\zeta)}
 &=-E+B_\alpha-T_{r,\alpha}\notag\\
 &=-[(3r+2)\alpha+3r^2+3r+1].\label{eq:oddloglimit}
\end{align}
Here $B_\alpha-T_{r,\alpha}=-(r+1)(\alpha+r)$. The relevant
determinants are nonzero along all sufficiently large members of this
sequence by the same endpoint estimates.

Subtracting \eqref{eq:evenloglimit} and \eqref{eq:oddloglimit} in
the appropriate order shows that the limit in \eqref{eq:Vlogder} is
$C_n(\lambda)$. Since
\[
 \deg H_n+\deg H_{n+1}
 =\left\lfloor\frac{n^2}4\right\rfloor
  +\left\lfloor\frac{(n+1)^2}4\right\rfloor=\frac{n(n+1)}2,
\]
boundedness of the ratio along an unbounded positive sequence proves
the degree upper bound, and its limit gives \eqref{eq:Vlead}.
The coefficients are polynomials in $\lambda$, so these statements
extend identically from $\lambda>-\tfrac12$. Since \eqref{eq:Vlead}
is not identically zero, the formal degree is exact.

For $\lambda>-\tfrac12$, every $\ell_n$ is nonzero.
The factor $C_{2r}$ is positive. The factor $C_1$ is one, and for
$r\ge1$ the unique zero of $C_{2r+1}$ is
$\lambda=\tfrac32-r+1/r$. Only $r=1,2$ put this zero in the
admissible parameter interval. The exact calculations below determine
the remaining specialised degrees.
\end{proof}

At $\lambda=3/2$, the defining recurrence gives
\begin{equation}\label{eq:exception3}
 V_3=-4608(13u^5-402u^4+4428u^3-18441u^2+14175u-3645).
\end{equation}
At $\lambda=0$, it gives
\begin{equation}\label{eq:exception5}
 \begin{split}
 V_5=181440u^9(&208u^5-9468u^4+154323u^3\\
              &-973368u^2+1369008u-648000).
 \end{split}
\end{equation}
Their leading coefficients are $-59904$ and $37739520$.
These are finite exact polynomial identities, not numerical evaluations
of Bessel functions. They can be checked directly from
\eqref{eq:Prec}, \eqref{eq:Hdef} and \eqref{eq:Htilde}.

The source writes its coefficients as functions of $\lambda$, so the
natural reading of its second degree conjecture is the formal generic
statement. Theorem~\ref{thm:Vdegree} confirms that statement and refines
it at exceptional specialisations; the exceptions are not presented as
a refutation of that generic interpretation.

\section{The recurrence-coefficient conjectures}\label{sec:recurrence}

Use the recurrence convention
\[
 \pi_{n+1}(t)=(t-\ii\alpha_n)\pi_n(t)-\beta_n\pi_{n-1}(t).
\]
At the Bessel-zero frequencies under consideration, its determinant
formulas are
\begin{equation}\label{eq:abdet}
 \alpha_n=\frac1\zeta\left(\frac{\wt H_n}{H_n}
                           -\frac{\wt H_{n+1}}{H_{n+1}}\right),
 \qquad \beta_n=-\frac{H_{n-1}H_{n+1}}{uH_n^2},\qquad u=\zeta^2.
\end{equation}
They hold as rational-function identities, and give actual recurrence
coefficients whenever the required determinants are nonzero. In the
source regime of rational $\lambda>-\tfrac12$ and nonzero zeros of
$J_{\lambda-1}$, nonvanishing follows from
\cite[Theorem~2.2]{MCM2008}.

\subsection{The general parameter}

Put $d_n=\lfloor n^2/4\rfloor$, $E_n=H_n/\ell_n$, and $E_0=E_1=1$.
The limits \eqref{eq:evenloglimit}--\eqref{eq:oddloglimit} show that
$\deg\wt H_n\le d_n$ and that $\wt H_n/H_n\to b_n$, where
\[
 b_{2r}=-2r(\lambda-\tfrac12+r),\qquad
 b_{2r+1}=-(3r+2)\lambda-3r^2-\tfrac32r.
\]
Set $S_n=n(2\lambda+n-1)$ and $e_n=S_n+b_n$. Thus
\begin{equation}\label{eq:general-e}
 e_{2r}=r(2\lambda+2r-1),\qquad
 e_{2r+1}=\frac r2(2\lambda+2r+1).
\end{equation}
For $n\ge2$, these are positive when $\lambda>-\tfrac12$.

\begin{lemma}\label{lem:constantshift}
The polynomial identity $\wt H_n(0)+S_nH_n(0)=0$ holds for every
$n\ge1$. Consequently
\begin{equation}\label{eq:Fdef}
 F_n(u)=\frac{\wt H_n(u)+S_nH_n(u)}{u\ell_ne_n},\qquad n\ge2,
\end{equation}
is monic of degree $d_n-1$.
\end{lemma}

\begin{proof}
At $u=0$, \eqref{eq:Prec} gives $P_j(0)=(-1)^j(b)_j$, $b=2\lambda$.
Factor $(b)_i$ from row $i$ and remove the row and column signs.
The unshifted determinant is then the evaluation determinant of the
monic polynomials $q_j(x)=(b+x)_j$, $0\le j<n$, at $x=0,\ldots,n-1$.
Its determinant is the corresponding Vandermonde. The shifted last
column uses $q_n$. Its coefficient of $x^{n-1}$ is
$nb+n(n-1)/2$. Reducing it modulo $\prod_{i=0}^{n-1}(x-i)$ adds a
further $n(n-1)/2$ to that coefficient. The ratio of the two evaluation
determinants is therefore $n(b+n-1)=S_n$. The extra column shift changes
the original sign by minus one. This proves the identity whenever
the factored quantities are nonzero, hence identically in $\lambda$.

The numerator in \eqref{eq:Fdef} is divisible by $u$. Its leading
coefficient is $\ell_n(S_n+b_n)=\ell_ne_n\ne0$, proving the degree
and monicity assertion.
\end{proof}

Define $X_{d_n}=E_n$ for $n\ge1$ and $X_{d_n-1}=F_n$ for $n\ge4$.
These definitions are consistent: the sequence $d_n$, $n\ge1$, is
strictly increasing, and $d_{n-1}<d_n-1<d_n$ for $n\ge4$. There are
therefore no collisions between numerator and denominator labels.
In particular $X_0=1$, $X_1=u-2\lambda$, and every assigned $X_j$
is monic of degree $j$; unused labels may be filled arbitrarily.

\begin{theorem}\label{thm:LZ25}
All four formulas of \cite[Conjecture~2.5]{LyuZhou2026} hold in their
stated ranges. Explicitly, with every $X$ evaluated at $u=\zeta^2$,
\begin{align}
 \alpha_{2k}={}&\frac{2(\lambda+2k)}\zeta
 +k\zeta\left[(2\lambda+2k-1)\frac{X_{k^2-1}}{X_{k^2}}
 -\frac{2\lambda+2k+1}{2}
       \frac{X_{k(k+1)-1}}{X_{k(k+1)}}\right],&&k\ge2,
 \label{eq:LZaeven}\\
 \alpha_{2k+1}={}&\frac{2(\lambda+2k+1)}\zeta
 +(2\lambda+2k+1)\zeta\left[
 \frac k2\frac{X_{k(k+1)-1}}{X_{k(k+1)}}
 -(k+1)\frac{X_{k(k+2)}}{X_{(k+1)^2}}\right],&&k\ge2,
 \label{eq:LZaodd}\\
 \beta_{2k}={}&-2k(2\lambda+2k-1)
       \frac{X_{k(k-1)}X_{k(k+1)}}{uX_{k^2}^2},&&k\ge1,
 \label{eq:betaeven}\\
 \beta_{2k+1}={}&\frac{X_{k^2}X_{(k+1)^2}}{uX_{k(k+1)}^2},&&k\ge0.
 \label{eq:betaodd}
\end{align}
\end{theorem}

\begin{proof}
By \eqref{eq:Fdef}, $\wt H_n/H_n=-S_n+ue_nF_n/E_n$.
Substitution in \eqref{eq:abdet}, using $S_{n+1}-S_n=2(\lambda+n)$,
gives
\[
 \alpha_n=\frac{2(\lambda+n)}\zeta+
       \zeta\left(e_n\frac{F_n}{E_n}-e_{n+1}\frac{F_{n+1}}{E_{n+1}}\right).
\]
The values \eqref{eq:general-e} give the first two formulas.
The ratios of the leading coefficients
\eqref{eq:Hevenlead}--\eqref{eq:Hoddlead} give the last two.
\end{proof}

The lower-index restriction on the $\alpha$ formulas is necessary for
this common-index notation. Directly,
$F_2=1$, $F_3=u-4\lambda^2-4\lambda$, and $E_2=u-2\lambda$.
The last two degree-one polynomials are generally different.
The source uses the displayed common-index formulas only for $k\ge2$;
its separate formulas for $\alpha_2,\alpha_3$ retain the appropriate
small-index numerator. For fixed admissible $\lambda$ the constructed
$X_j$ are ordinary polynomials in $u$; formally their coefficients lie
in $\Q(\lambda)$.

\subsection{The exceptional parameter \texorpdfstring{$\lambda=0$}{lambda=0}}

At $\lambda=0$, powers of $u$ cancel from the Hankel determinants,
leading to the different index pattern of
\cite[Conjecture~2.7]{LyuZhou2026}. We now establish that pattern
and all its recurrence formulas.

\begin{lemma}\label{lem:valuation}
For $n\ge1$, the exact valuation of $H_n(u;0)$ at $u=0$ is $n-1$.
Writing $m=n-1$, its first coefficient is
\begin{equation}\label{eq:lowcoefficient}
 [u^{n-1}]H_n(u;0)=
 \left(-\frac12\right)^m
 \left(1+\frac{m(m+1)}2\right)\prod_{j=0}^{m-1}j!(j+1)!.
\end{equation}
If $H_n=u^{n-1}R_n$ and $\wt H_n=u^{n-1}\wt R_n$, then
\begin{equation}\label{eq:zeroRshift}
 \frac{\wt R_n(0)}{R_n(0)}=-A_n,
 \qquad A_n=\frac{[n(n-1)]^2}{n(n-1)+2}.
\end{equation}
\end{lemma}

\begin{proof}
The case $n=1$ is immediate, so assume $n\ge2$. Every entry in rows
$1,\ldots,n-1$ of $H_n$ and of $\wt H_n$ has positive moment index, hence
is divisible by $u$, because every $P_j$ with $j\ge1$ is (at $\lambda=0$,
$P_1=0$ and \eqref{eq:Prec} gives $u\mid P_{k+2}$ inductively); factoring $u$ from those rows proves
$u^{n-1}\mid H_n$ and $u^{n-1}\mid\wt H_n$, so $R_n$ and $\wt R_n$ are
polynomials.
Every $P_j$, $j\ge1$, is divisible by $u$. Its first coefficient is
$q_1=0$, $q_2=-1$, and
$q_j=(-1)^{j-1}(j-1)!/2$ for $j\ge3$, by \eqref{eq:Prec}.
Factor $u$ from rows $1,\ldots,n-1$ in $H_n$ and take the constant
term of the remaining determinant. The first row then reduces to
its top entry, giving
\[
 [u^{n-1}]H_n=
 \left(-\frac12\right)^m
 \det[\Gamma(i+j+2)+\delta_{i0}\delta_{j0}]_{i,j=0}^{m-1}.
\]
The last matrix is the positive Gram matrix for
\begin{equation}\label{eq:Laguerreatom}
                         \dd\rho(t)=te^{-t}\dd t+\delta_0,
                         \qquad t\in[0,\infty),
\end{equation}
so it is nonsingular.

For completeness its determinant and shifted ratio can be computed
together. Let $p_j(t)=(-1)^jj!L_j^{(1)}(t)$ be monic Laguerre
polynomials for $te^{-t}\dd t$. Their squared norms are
$h_j=j!(j+1)!$, their values at zero are $(-1)^j(j+1)!$, and the
coefficient of $t^{m-1}$ in $p_m$ is $-m(m+1)$. For the kernel
$K_{m-1}(t,s)=\sum_{j=0}^{m-1}p_j(t)p_j(s)/h_j$,
\[
                         K_{m-1}(0,0)=\frac{m(m+1)}2.
\]
The matrix determinant lemma gives the factor
$1+K_{m-1}(0,0)$ multiplying the Laguerre determinant, proving
\eqref{eq:lowcoefficient}. The monic polynomial orthogonal for
\eqref{eq:Laguerreatom} is
\[
 Q_m(t)=p_m(t)-\frac{p_m(0)}{1+K_{m-1}(0,0)}K_{m-1}(t,0).
\]
Its orthogonality follows directly from the reproducing property of
the kernel. Put $B=m(m+1)$. Its subleading coefficient is
$-B+B/(1+B/2)=-B^2/(B+2)$.
The shifted-column determinant formula for that coefficient,
together with the extra minus sign from the shifted column in the
original $q_j$, identifies it with $\wt R_n(0)/R_n(0)$.
This proves \eqref{eq:zeroRshift}.
\end{proof}

Put
\begin{equation}\label{eq:sndef}
 s_n=\deg R_n=d_n-n+1=
                  \left\lfloor\frac{(n-2)^2}{4}\right\rfloor.
\end{equation}
At $\lambda=0$, the limits of the shifted ratios are
\[
 b_n=\begin{cases}
       -n(n-1)/2,&n\text{ even},\\
       -3n(n-1)/4,&n\text{ odd}.
      \end{cases}
\]
Set $\eta_n=A_n+b_n$. Explicitly,
\begin{equation}\label{eq:eta}
 \eta_{2k}=\frac{k(k-1)(2k-1)(2k+1)}{k(2k-1)+1},\qquad
 \eta_{2k+1}=\frac{k(k-1)(2k+1)(2k+3)}{2[k(2k+1)+1]}.
\end{equation}
These are nonzero for $n\ge4$. By \eqref{eq:zeroRshift},
\[
 T_n(u)=\frac{\wt R_n(u)+A_nR_n(u)}u
\]
is a polynomial of degree $s_n-1$ with leading coefficient
$\ell_n\eta_n$. Define
\[
 Y_{s_n}=R_n/\ell_n\quad(n\ge2),\qquad
 Y_{s_n-1}=T_n/(\ell_n\eta_n)\quad(n\ge6).
\]
The two definitions of $Y_0$ from $n=2,3$ agree and equal one;
$Y_1=u-21/4$. For $n\ge6$,
$s_{n-1}<s_n-1<s_n$, so all further numerator labels are distinct
and do not collide with denominator labels. Each assigned $Y_j$ is
monic of degree $j$.

\begin{theorem}\label{thm:LZ27}
All formulas of \cite[Conjecture~2.7]{LyuZhou2026} hold in their
stated ranges. A unified form for its diagonal formulas is
\begin{equation}\label{eq:LZzeroalpha}
 \begin{split}
 \alpha_n={}&
 \frac{2n^3(n^2+3)}{(n^2-n+2)(n^2+n+2)\zeta}\\
 &+\zeta\left(
 \eta_n\frac{Y_{s_n-1}(u)}{Y_{s_n}(u)}
 -\eta_{n+1}\frac{Y_{s_{n+1}-1}(u)}{Y_{s_{n+1}}(u)}\right),
 \qquad n\ge6.
 \end{split}
\end{equation}
The two off-diagonal formulas are
\begin{align}
 \beta_{2k}&=-\frac{2k(2k-1)}u
 \frac{Y_{(k-1)(k-2)}Y_{k(k-1)}}{Y_{(k-1)^2}^2},&&k\ge1,
 \label{eq:LZzerobeven}\\
 \beta_{2k+1}&=\frac1u
 \frac{Y_{(k-1)^2}Y_{k^2}}{Y_{k(k-1)}^2},&&k\ge1.
 \label{eq:LZzerobodd}
\end{align}
\end{theorem}

\begin{proof}
The construction gives
$\wt H_n/H_n=-A_n+u\eta_nY_{s_n-1}/Y_{s_n}$ for $n\ge6$.
Insert this in \eqref{eq:abdet} and use
\[
 A_{n+1}-A_n=\frac{2n^3(n^2+3)}{(n^2-n+2)(n^2+n+2)}
\]
to obtain \eqref{eq:LZzeroalpha}. Inserting
$H_n=u^{n-1}\ell_nY_{s_n}$ in the second formula of
\eqref{eq:abdet}, and using
$\ell_{2k-1}\ell_{2k+1}/\ell_{2k}^2=2k(2k-1)$ and
$\ell_{2k}\ell_{2k+2}/\ell_{2k+1}^2=-1$, gives the two $\beta$
formulas. Substituting $n=2k$, $n=4j-1$, and $n=4j+1$ into
\eqref{eq:LZzeroalpha} gives the three diagonal displays in the
source, respectively for $k\ge3$, $j\ge2$, and $j\ge2$.
\end{proof}

The source's exclusions of the diagonal indices four and five matter.
At $n=5$, the monic secondary numerator is $u-177/8$, whereas the
already fixed denominator polynomial $Y_1$ is $u-21/4$.
The stated ranges avoid identifying those different polynomials.
At $n=4$ the secondary numerator is simply $Y_0=1$.

\subsection{Comparison with the polynomials listed by Lyu and Zhou}
Lyu and Zhou list the polynomials $X_j$ and $Y_j$ only at low indices
\cite[Appendix~C]{LyuZhou2026}. Table~\ref{tab:LZ} records an exact
symbolic comparison of every listed polynomial with the monic
constructions used above ($X$ for generic $\lambda$, $Y$ at $\lambda=0$).
All eighteen agree exactly, as do the base cases $X_0=1$,
$X_1=u-2\lambda$ and $Y_1=u-21/4$.

\begin{table}[ht]
\centering
\begin{tabular}{llll}
\hline
Family & Index $j$ & Manuscript polynomial & Exact agreement\\
\hline
$X_{2}$ & 2 & $H_{3}/\ell_{3}$ (monic) & yes\\
$X_{3}$ & 3 & $F_{4}$ (monic) & yes\\
$X_{4}$ & 4 & $H_{4}/\ell_{4}$ (monic) & yes\\
$X_{5}$ & 5 & $F_{5}$ (monic) & yes\\
$X_{6}$ & 6 & $H_{5}/\ell_{5}$ (monic) & yes\\
$X_{8}$ & 8 & $F_{6}$ (monic) & yes\\
$X_{9}$ & 9 & $H_{6}/\ell_{6}$ (monic) & yes\\
$Y_{2}$ & 2 & $R_{5}/\ell_{5}$ (monic) & yes\\
$Y_{3}$ & 3 & $T_{6}$ (monic) & yes\\
$Y_{4}$ & 4 & $R_{6}/\ell_{6}$ (monic) & yes\\
$Y_{5}$ & 5 & $T_{7}$ (monic) & yes\\
$Y_{6}$ & 6 & $R_{7}/\ell_{7}$ (monic) & yes\\
$Y_{8}$ & 8 & $T_{8}$ (monic) & yes\\
$Y_{9}$ & 9 & $R_{8}/\ell_{8}$ (monic) & yes\\
$Y_{11}$ & 11 & $T_{9}$ (monic) & yes\\
$Y_{12}$ & 12 & $R_{9}/\ell_{9}$ (monic) & yes\\
$Y_{15}$ & 15 & $T_{10}$ (monic) & yes\\
$Y_{16}$ & 16 & $R_{10}/\ell_{10}$ (monic) & yes\\
\hline
\end{tabular}
\caption{Exact agreement of the manuscript's polynomials with those listed in \cite[Appendix~C]{LyuZhou2026}.}\label{tab:LZ}
\end{table}

\section{Consequences and limits}\label{sec:consequences}

\subsection{Rationality in special-frequency existence}

The rationality assumption in the all-degree existence theorem of
\cite[Theorem~2.2]{MCM2008} cannot be removed throughout the real
parameter interval.

\begin{proposition}\label{prop:irrational}
For every integer $m\ge2$, there is an irrational $\lambda_m>1/2$ such
that, with $\zeta_m=j_{\lambda_m-1,m}$,
\[
                   D_3^{(\lambda_m-1/2)}(\zeta_m)=0.
\]
There are infinitely many distinct such parameters. At each of these
points the monic orthogonal polynomial of degree three does not exist,
although every even degree exists uniquely.
\end{proposition}

\begin{proof}
The recurrence \eqref{eq:Prec} gives
\[
 H_3(u;\lambda)=2(2\lambda+1)
        [u^2-\lambda(2\lambda+5)u+4\lambda^2].
\]
For $\lambda\ge1/2$, set
\[
 F_m(\lambda)=j_{\lambda-1,m}^4
 -\lambda(2\lambda+5)j_{\lambda-1,m}^2+4\lambda^2.
\]
Positive Bessel zeros vary continuously with the order in this range.
At $\lambda=1/2$, put $u_m=(m-1/2)^2\pi^2$. Then
$F_m(1/2)=u_m^2-3u_m+1>0$ for $m\ge2$. The fixed-index large-order
expansion $j_{\nu,m}=\nu+O(\nu^{1/3})$
\cite[\S10.21, Eq.~10.21.32]{DLMF} gives
$F_m(\lambda)/\lambda^4\to-1$. The intermediate value theorem
therefore gives a root $\lambda_m>1/2$.

At a positive zero of $J_{\lambda-1}$, the function $J_\lambda$ is
nonzero: a common zero would also make the derivative of
$J_{\lambda-1}$ vanish at an ordinary point of its differential
equation. Thus \eqref{eq:DH} converts $F_m(\lambda_m)=0$ into the
asserted determinant zero. If $\lambda_m$ were rational,
\cite[Theorem~2.2]{MCM2008} would give nonvanishing of every
determinant there. Hence $\lambda_m$ is irrational. A fixed parameter
permits at most two positive $u$ solving the displayed quadratic, so
finitely many parameters cannot account for all indices $m\ge2$.

Finally, let $p_2$ be the monic quadratic, which exists because
$D_2>0$. At $D_3=0$, its norm is zero, so it is orthogonal through
degree two. If a monic orthogonal cubic existed, symmetry of the
bilinear form would also make $p_2$ orthogonal to degree three. The
size-four moment matrix would then be singular, contrary to $D_4>0$.
The even-degree conclusion follows from Theorem~\ref{thm:main}.
\end{proof}

This excludes a natural unrestricted-real-parameter extension of the
2008 existence theorem; no published conjecture of that extension is
being attributed to its authors.

\subsection{Symmetry is insufficient}

Consider the symmetric positive probability measure
\[
 \nu=\frac16(\delta_{-1}+\delta_{-1/2}+2\delta_0+
                         \delta_{1/2}+\delta_1).
\]
Form its Fourier Hankel determinants using
$\mu_j=\int t^je^{\ii xt}\dd\nu(t)$. At $x=2\pi$, odd moments
vanish and
\[
       \mu_0=\frac13,\qquad \mu_2=\frac14,
       \qquad\mu_4=\frac5{16},\qquad\mu_6=\frac{21}{64}.
\]
Thus $D_2=1/12>0$, whereas parity-block decomposition gives
\begin{equation}\label{eq:atomic}
 D_4=\left(\frac13\frac5{16}-\frac1{16}\right)
       \left(\frac14\frac{21}{64}-\frac{25}{256}\right)
     =-\frac1{1536}<0.
\end{equation}
The Gegenbauer structure in the positivity theorem is substantive.
Positive-definite functions may take negative values, so this example
does not conflict with the Fourier-Wronskian positive-definiteness
results of \cite{DimitrovXu2019}.

\subsection{Spherical-Bessel determinants and critical length}

Define
\begin{equation}\label{eq:fW}
 f_n(x)=\sqrt{\frac\pi2}x^{n+1/2}J_{n+1/2}(x),\qquad
 W_k[f](x)=(-1)^{k(k-1)/2}
                      \det[f^{(i+j)}(x)]_{i,j=0}^{k-1}.
\end{equation}
Appendix~\ref{app:duality} proves that, for $1\le k\le n+1$,
\begin{equation}\label{eq:dualitymain}
 W_k[f_n](x)=
 \frac{2^{k(k-n-2)}(n!)^k}
      {\displaystyle\prod_{q=n+1-k}^{n}(q!)^2}
 x^{k(2n+2-k)}D_k^{(n+1-k)}(x).
\end{equation}
All constants are positive. Theorem~\ref{thm:main} therefore gives
\begin{corollary}\label{cor:evenW}
For every $n\ge0$, every even $k$ with $2\le k\le n+1$, and every
$x>0$, one has $W_k[f_n](x)>0$.
\end{corollary}

Let
\[
 U_n=\spanop\{x^j\cos x,x^j\sin x:0\le j\le n\}.
\]
Its critical length is the supremum of interval lengths on which every
nonzero member has fewer than $2n+2$ zeros, counted with multiplicity.
Carnicer, Mainar and Pe\~na \cite{CMP2023} conjecture that this length
is the first positive zero $j_{n+1/2,1}$ of $J_{n+1/2}$.

\begin{corollary}\label{cor:CMP3}
The critical length of $U_3$ is $j_{7/2,1}$.
\end{corollary}

\begin{proof}
The space $U_n$ is invariant under translation and reflection, and
$u_j=f_n^{(2n+1-j)}$, $0\le j\le2n+1$, is a canonical basis at zero.
Here canonical means $u_j^{(i)}(0)=0$ for $i<j$ and
$u_j^{(j)}(0)\ne0$, as in \cite[Eq.~(2)]{CMP2023}.
Indeed $f_n$ has a zero of order $2n+1$ there and belongs to
$\ker(D^2+1)^{n+1}$. The critical length is finite: $f_n$ has
$2n+1$ zeros at the origin, counted with multiplicity, and another at
its first positive zero. Proposition~3 of \cite{CMP2023} states that the
critical length is the least positive zero among the suffix
Wronskians with $j>(2n+1)/2$. Such a suffix has size
$k=2n+2-j$, so precisely $1\le k\le n+1$ occurs. Reversing its
columns identifies its Wronskian with $W_k[f_n]$.

For $n=3$, only $k=1,2,3,4$ is required. The function $W_1=f_3$ is
positive on $(0,j_{7/2,1})$ and vanishes at its right endpoint.
Corollary~\ref{cor:evenW} gives positivity of $W_2,W_4$ for every
$x>0$. Kounchev and Render \cite[Corollary~21]{KounchevRender2025}
(an unrefereed preprint, on which this corollary depends) prove that $W_3[f_3]$ is positive on $(0,j_{7/2,1})$. Thus none of
the required suffix Wronskians vanishes earlier, and one vanishes at
the stated endpoint.
\end{proof}

For general $n$, \eqref{eq:dualitymain} leaves a first-zero comparison
for odd determinants of size at least five. Infinite oscillation and
simple zeros do not supply that comparison. No resolution of the
general critical-length conjecture is claimed.

\section{Checks and the scope of the literature comparison}\label{sec:scope}

At $\alpha=0$ and size two,
\[
 D_2^{(0)}(x)=\frac{4(x^2-\sin^2x)}{x^4},\qquad D_2^{(0)}(0)=\frac43.
\]
This checks the sign and $C_{1,0}=4$ in \eqref{eq:eventail}. The general
coefficient formula gives $C_{2,0}=256$, consistent with the exponent-zero endpoint
calculation in \cite{Celsus2022}. At odd size one,
$D_1^{(0)}(x)=2\sin x/x$ and $T_{0,0}=0$.
Exact symbolic calculations from \eqref{eq:Prec} for $n=1,\ldots,5$
give degrees $0,1,2,4,6$ and leading coefficients
\[
 1,\quad -1,\quad 2(2\lambda+1),\quad4(2\lambda+1)^2,
 \quad32(2\lambda+1)^3(2\lambda+3).
\]
They verify normalisation and signs; the all-rank arguments above are
independent of these finite checks.

The companion script \texttt{verify.py} records its results in
\texttt{verification\_results.json}. At 90 decimal digits it checked
120 even determinants: sizes $2,4,6,8$, frequencies
$0,0.1,\pi,10,50$, and exponents
$-0.99,-0.5,0,0.5,1.4142135623730950488,3$.
All computed determinants were positive. The moments were cross-checked
against an independent hypergeometric formula. The same script checks
\eqref{eq:selbergzero}, the exact $D_2^{(0)}$ identity, the tail constants,
and the Toda identity. Its exact symbolic checks cover $H_n$ through
$n=5$, generic shifted numerators through $n=4$, the general-parameter
residual polynomials through $n=5$, and the zero-parameter reduced
determinants and shifted ratios through $n=9$. It also verifies both
exceptional polynomials \eqref{eq:exception3}--\eqref{eq:exception5} and the atomic
example \eqref{eq:atomic}. These computations are checks of formulas
and normalisations, not a numerical proof of the theorems.

The comparison made for this draft includes the full 2008 paper, the
2022 exponent-zero existence proof, Dimitrov--Xu's Fourier-Wronskian
results, the 2025 rank-three critical-length result, and Lyu--Zhou's
2026 recurrence paper. The finite verification record
\cite{Ideosphere2026} reports degree checks through $n=16$ and more
limited recurrence-coefficient checks. Such a finite verification is
different in scope from the all-rank arguments here.

A close classical comparison is also important. The Karlin--Szeg\H{o}
even-determinant inequalities, as set out in
\cite[Section~2.2]{BariczPogany2014}, differentiate
$u^{-\nu/2}J_\nu(2\sqrt u)$ with respect to $u$ and hence lead to
order-shifted Bessel entries. Our determinant differentiates
$x^{-\nu}J_\nu(x)$ with respect to $x$. A nonlinear substitution does
not generally preserve derivative Hankel determinants: writing
$\mathcal H_2[f]=ff''-(f')^2$, the chain rule gives
\[
 \mathcal H_2[f\circ\phi]
 =\phi'^2(\mathcal H_2[f]\circ\phi)
   +\phi''((ff')\circ\phi).
\]
Thus that classical result is not an immediate reparametrisation of
Theorem~\ref{thm:main}. This comparison rules out the direct
identification; it does not claim to exclude every deeper connection.

\paragraph{The earlier claim.}
Ashfaque \cite{Ashfaque2026} (ResearchGate, 24 July 2026) states in its
abstract and Proposition~3.1 that the symbolic recurrence formulas
conjectured by Lyu and Zhou follow from a discrete Painlev\'e~IV and
tau-function identification. Our description is based on a full reading
of that six-page document on 2 November 2026; the host later blocked
automated access, so readers should consult the original. In its
notation the weight is $(1-x^2)^{\lambda-1/2}e^{\ii\zeta x}$ with
$J_{\lambda-1}(\zeta)=0$, the recurrence is
$xP_n=P_{n+1}+\ii\alpha_nP_n+\beta_nP_{n-1}$ with real $\alpha_n$, and
it sets $u_n=\zeta\alpha_n-(2n+2\lambda-1)$. Proposition~3.1 defines
$D_n^{(1)}=\det[\mu_{i+j+1}]_{i,j=0}^{n-1}$, shifting every moment
index, and asserts
\begin{equation}\label{eq:ashfaqueA}
 u_n=\zeta\frac{D_n^{(1)}}{D_n}-(2n+2\lambda-1).
\end{equation}
This cannot hold as printed. For real $\zeta$,
$\overline{\mu_j}=(-1)^j\mu_j$, so the argument used in
Section~\ref{sec:continuation} shows that $D_n^{(1)}$ is purely imaginary
when $n$ is odd, while $D_n$ is real. At $n=1$,
$\mu_1/\mu_0=\ii J_{\lambda+1}(\zeta)/J_\lambda(\zeta)=2\ii\lambda/\zeta$ at
every nonzero real zero of $J_{\lambda-1}$ (where $J_\lambda(\zeta)\ne0$,
since adjacent Bessel functions have no common nonzero zero), by the three-term Bessel recurrence. The
right side of \eqref{eq:ashfaqueA} is then $2\ii\lambda-(2\lambda+1)$,
which is not real when $\lambda\ne0$, whereas the left side is real. At
$\lambda=0$ the two sides are $-1$ and $0$ respectively, using
Lyu and Zhou's initial value $\alpha_1=1/\zeta$. The failure is not
an index slip: the right side alternates between imaginary and real as
the parity of $n$ changes, so no fixed index shift or constant phase
repairs it for all $n$. Moreover the subleading coefficient of a monic orthogonal
polynomial comes from a single last-column shift, and $\alpha_n$ is a
difference of two such ratios. A corrected argument along those lines
may exist; we make no judgement about one. The arguments in this
manuscript do not use that proposition or the claimed resolution, but
the earlier public claim is part of the priority record and we do not
claim to be first to assert these formulas.

This is a bounded comparison, not an exhaustive priority determination.
In particular, no claim of historical firstness is inferred solely
from the absence of a located valid proof. The central mathematical
claims of the draft are the extension in Theorem~\ref{thm:main}, the
explicit all-rank algebraic results in
Theorems~\ref{thm:Hdegree} and \ref{thm:Vdegree}, and the stated
critical-length corollary. They are provided with proofs so that their
validity and their relation to earlier work can be assessed separately.

\appendix
\section{An exact Bessel--Jacobi determinant identity}\label{app:duality}

This appendix proves \eqref{eq:dualitymain}, including its constant.
The two standard inputs are the normalised unitary HCIZ integral
\cite[Eq.~(2)]{McSwiggen2019} and the Jacobi law for random projection
compressions \cite[Theorem~2.2]{Collins2005}. The determinant reduction
between them is given explicitly.

\subsection{A confluent orbital-integral identity}

Let $1\le k\le N$ and initially let $a_1,\ldots,a_N$ be distinct
complex numbers. Put
\[
 p(z)=\prod_{j=1}^N(z-a_j),\qquad
 g(x)=\sum_{j=1}^N\frac{e^{a_jx}}{p'(a_j)},\qquad
 A=\diag(a_1,\ldots,a_N),\quad B=\diag(0_{N-k},1_k).
\]
Normalise Haar measure on $U(N)$ to mass one and set
\[
 I_{N,k}(A;x)=\int_{U(N)}e^{x\tr(AUBU^*)}\dd U,\qquad
 Q_{N,k}=\frac{\prod_{j=N-k}^{N-1}j!}{\prod_{j=0}^{k-1}j!}.
\]
We claim that
\begin{equation}\label{eq:HCIZW}
            I_{N,k}(A;x)=Q_{N,k}x^{-k(N-k)}W_k[g](x).
\end{equation}
Start with the normalised HCIZ formula
\begin{equation}\label{eq:HCIZ}
 \int_{U(N)}e^{x\tr(AUCU^*)}\dd U
 =\left(\prod_{j=1}^{N-1}j!\right)
 \frac{\det[e^{xa_ic_j}]_{i,j=1}^N}
 {x^{N(N-1)/2}\Delta(a)\Delta(c)}.
\end{equation}
Analyticity extends the usual real diagonal formula to complex
diagonal entries. Coalesce $N-k$ eigenvalues of $C$ at zero and $k$
at one. The confluent columns are
\[
 \frac{(xa_i)^j}{j!}\quad(0\le j<N-k),\qquad
 \frac{(xa_i)^\ell e^{xa_i}}{\ell!}\quad(0\le\ell<k).
\]
The cross-cluster Vandermonde factors are one. Factoring powers and
factorials gives
\begin{equation}\label{eq:Mdet}
 I_{N,k}(A;x)=Q_{N,k}x^{-k(N-k)}\frac{\det M(x)}{\Delta(a)},
 \quad
 M=[a_i^j\ (0\le j<N-k)\mid a_i^\ell e^{xa_i}\ (0\le\ell<k)].
\end{equation}

Let $V=[a_i^j]_{1\le i\le N,\;0\le j<N}$. The first $N-k$ columns
of $V^{-1}M$ are coordinate vectors, so its determinant is that of its
bottom-right $k$-by-$k$ block. In a column labelled by $\ell$, its
entries are coefficients of the polynomial interpolating
$a_i^\ell e^{xa_i}$ at the nodes. The coefficient of $z^{N-1-s}$ in
the $i$th Lagrange polynomial is
\[
 \frac{(-1)^se_s(a_1,\ldots,\widehat{a_i},\ldots,a_N)}{p'(a_i)}
 =\frac{a_i^s+\sum_{q=0}^{s-1}c_{s,q}a_i^q}{p'(a_i)},
\]
where the $c_{s,q}$ depend only on the full node set. The corresponding
entry is therefore
$g^{(s+\ell)}+\sum_{q<s}c_{s,q}g^{(q+\ell)}$.
The bottom rows have $s=k-1,\ldots,0$. Reverse them, producing
$(-1)^{k(k-1)/2}$, and eliminate the lower triangular terms by unit
triangular row operations. The resulting determinant is exactly
$W_k[g]$. This proves \eqref{eq:HCIZW}.

The formulas extend continuously to coalescing nodes. For example,
$g$ is the exponential divided difference, equivalently
\[
                         g(x)=\frac1{2\pi\ii}\oint\frac{e^{xz}}{p(z)}\dd z,
\]
where the contour encloses all roots of $p$.

\subsection{Coalescing at the spherical-Bessel spectrum}

Take $N=2n+2$, $m=n+1$, and $A=\diag(-\ii I_m,\ii I_m)$.
Then $p(z)=(z^2+1)^{n+1}$ and $g$ is the fundamental solution of
\[
 (D^2+1)^{n+1}g=0,\qquad
 g^{(j)}(0)=0\ (0\le j<2n+1),\qquad g^{(2n+1)}(0)=1.
\]
The spherical-Bessel function $f_n$ belongs to this solution space
and satisfies
\[
 f_n(x)=\frac{x^{2n+1}}{(2n+1)!!}+O(x^{2n+3}),\qquad
 \frac{(2n+1)!}{(2n+1)!!}=2^nn!.
\]
Hence $g=f_n/(2^nn!)$, and \eqref{eq:HCIZW} becomes
\begin{equation}\label{eq:orbitalW}
 I_{N,k}(A;x)=\frac{Q_{N,k}}{(2^nn!)^k}
                     x^{-k(N-k)}W_k[f_n](x).
\end{equation}

Write $A=\ii(2P-I_N)$, where $P$ is a rank-$m$ orthogonal projection.
If $V$ consists of the $k$ columns of $U$ selected by $B$, then
$\tr(AUBU^*)=\ii(2\tr(V^*PV)-k)$. For $k\le m$, the eigenvalues
$s_1,\ldots,s_k$ of $V^*PV$ have Jacobi density
\begin{equation}\label{eq:Jacobilaw}
 \text{constant}\cdot\Delta(s)^2
        \prod_{j=1}^ks_j^{m-k}(1-s_j)^{m-k},\qquad 0<s_j<1.
\end{equation}
After $t_j=2s_j-1$, this is proportional to
$\Delta(t)^2\prod(1-t_j^2)^{m-k}$. Andr\'eief's identity therefore
identifies the same orbital integral as
\begin{equation}\label{eq:orbitalD}
                  I_{N,k}(A;x)=\frac{D_k^{(n+1-k)}(x)}{D_k^{(n+1-k)}(0)}.
\end{equation}

Finally, \eqref{eq:selbergzero} gives
\[
 Q_{2n+2,k}D_k^{(n+1-k)}(0)
       =2^{k(2n+2-k)}\prod_{q=n+1-k}^{n}(q!)^2.
\]
Combining \eqref{eq:orbitalW} and \eqref{eq:orbitalD} proves
\eqref{eq:dualitymain}. Both sides are entire, so the identity extends
to complex frequency, including zero by continuation.

For $k=1$ the identity reduces to
$f_n=x^{2n+1}D_1^{(n)}/(2^{n+1}n!)$. For $n=1,k=2$, it gives
$W_2[f_1]=x^2-\sin^2x$, agreeing with
$D_2^{(0)}=4(x^2-\sin^2x)/x^4$ and the constant $1/4$.

\begin{thebibliography}{99}

\bibitem{Celsus2022}
A.~F. Celsus, A. Dea\~no, D. Huybrechs and A. Iserles,
\emph{The kissing polynomials and their Hankel determinants},
Transactions of Mathematics and Its Applications \textbf{6} (2022), tnab005.
\href{https://doi.org/10.1093/imatrm/tnab005}{doi:10.1093/imatrm/tnab005}.

\bibitem{MCM2008}
G.~V. Milovanovi\'c, A.~S. Cvetkovi\'c and Z.~M. Marjanovi\'c,
\emph{Orthogonal polynomials for the oscillatory-Gegenbauer weight},
Publications de l'Institut Math\'ematique, Nouvelle s\'erie
\textbf{84}(98) (2008), 49--60.
\href{https://doi.org/10.2298/PIM0898049M}{doi:10.2298/PIM0898049M}.

\bibitem{DimitrovXu2019}
D.~K. Dimitrov and Y. Xu,
\emph{Wronskians of Fourier and Laplace transforms},
Transactions of the American Mathematical Society \textbf{372} (2019),
4107--4125.
\href{https://doi.org/10.1090/tran/7809}{doi:10.1090/tran/7809}.

\bibitem{LyuZhou2026}
S. Lyu and X. Zhou,
\emph{Recurrence coefficients of the orthogonal polynomials for
oscillatory Jacobi-type weight functions}, preprint (2026).
\href{https://arxiv.org/abs/2607.19797v1}{arXiv:2607.19797v1}.

\bibitem{ADHW2014}
Andreas Asheim, Alfredo Deaño, Daan Huybrechs and Haiyong Wang,
\emph{A Gaussian quadrature rule for oscillatory integrals on a bounded interval},
Discrete and Continuous Dynamical Systems \textbf{34}(3) (2014), 883--901.
\href{https://doi.org/10.3934/dcds.2014.34.883}{doi:10.3934/dcds.2014.34.883}.

\bibitem{DeanoHuybrechs2009}
Alfredo Deaño and Daan Huybrechs,
\emph{Complex Gaussian quadrature of oscillatory integrals},
Numerische Mathematik \textbf{112}(2) (2009), 197--219.
\href{https://doi.org/10.1007/s00211-008-0209-z}{doi:10.1007/s00211-008-0209-z}.

\bibitem{DHI2017}
Alfredo Deaño, Daan Huybrechs and Arieh Iserles,
\emph{Computing Highly Oscillatory Integrals},
Society for Industrial and Applied Mathematics, 2017.
\href{https://doi.org/10.1137/1.9781611975123}{doi:10.1137/1.9781611975123}.

\bibitem{Barhoumi2021}
A.~B. Barhoumi,
\emph{Strong asymptotics of Jacobi-type kissing polynomials},
Integral Transforms and Special Functions \textbf{32}(5--8) (2021),
377--394.
\href{https://doi.org/10.1080/10652469.2021.1923707}
{doi:10.1080/10652469.2021.1923707}.

\bibitem{BariczPogany2014}
\'A. Baricz and T.~K. Pog\'any,
\emph{Tur\'an determinants of Bessel functions},
Forum Mathematicum \textbf{26}(1) (2014), 295--322.
\href{https://doi.org/10.1515/form.2011.160}{doi:10.1515/form.2011.160};
consulted author preprint
\href{https://arxiv.org/abs/1101.4624v2}{arXiv:1101.4624v2}.

\bibitem{CMP2023}
J.~M. Carnicer, E. Mainar and J.~M. Pe\~na,
\emph{Spherical Bessel functions and critical lengths},
Revista de la Real Academia de Ciencias Exactas, F\'isicas y Naturales.
Serie A. Matem\'aticas \textbf{117} (2023), article 29;
published online 27 November 2022.
\href{https://doi.org/10.1007/s13398-022-01365-8}
{doi:10.1007/s13398-022-01365-8}.

\bibitem{KounchevRender2025}
O. Kounchev and H. Render,
\emph{On the critical length conjecture for spherical Bessel functions in
CAGD}, preprint (2025).
\href{https://arxiv.org/abs/2505.09964v1}{arXiv:2505.09964v1}.

\bibitem{McSwiggen2019}
C. McSwiggen,
\emph{A new proof of Harish-Chandra's integral formula},
Communications in Mathematical Physics \textbf{365} (2019), 239--253.
\href{https://doi.org/10.1007/s00220-018-3259-9}
{doi:10.1007/s00220-018-3259-9}.

\bibitem{Collins2005}
B. Collins,
\emph{Product of random projections, Jacobi ensembles and universality
problems arising from free probability},
Probability Theory and Related Fields \textbf{133} (2005), 315--344.
\href{https://doi.org/10.1007/s00440-005-0428-5}
{doi:10.1007/s00440-005-0428-5}.

\bibitem{DLMF}
NIST Digital Library of Mathematical Functions,
\emph{\S10.21: Zeros of Bessel and Hankel functions},
especially \S10.21(iv) and Eq.~10.21.32, accessed 2 November 2026.
\url{https://dlmf.nist.gov/10.21}.

\bibitem{Ideosphere2026}
Ideosphere, \emph{Oscillatory Jacobi recurrence record},
finite-verification snapshot dated 7 September 2026, accessed
2 November 2026.
\url{https://ideosphere.ai/papers/oscjacobi-rec}.

\bibitem{Ashfaque2026}
J.~M. Ashfaque,
\emph{On the discrete Painlev\'e IV structure of oscillatory Gegenbauer
recurrence coefficients}, preprint dated 24 July 2026, accessed
2 November 2026.
\href{https://www.researchgate.net/publication/410769119_On_the_Discrete_Painleve_IV_Structure_of_Oscillatory_Gegenbauer_Recurrence_Coefficients}
{\nolinkurl{https://www.researchgate.net/publication/410769119}}.

\end{thebibliography}
\end{document}
